\CatchFileBetweenTags{\flavorTotalChiVal}{calculations/flavor_geometry.tex}{flavorTotalChiVal}
\CatchFileBetweenTags{\flavorReducedChiVal}{calculations/flavor_geometry.tex}{flavorReducedChiVal}
\CatchFileBetweenTags{\AtmosAngleVal}{calculations/flavor_geometry.tex}{AtmosAngleVal}

\section{The Zero-DOF Global Flavor Fit}

To rigorously test the validity of the structural duality framework, we performed a global $\chi^2$ consistency check across the complete flavor sector. This audit incorporates $17$ fundamental observables: all $9$ CKM matrix elements, the Jarlskog invariant, the Unitarity Triangle angles ($\delta_q, \gamma, \sin2\beta$), and the $4$ PMNS matrix parameters ($\theta_{12}, \theta_{23}, \theta_{13}, \delta_{CP}^\nu$).

Crucially, standard global fits (like those performed by the Particle Data Group or NuFIT) achieve their precision by allowing underlying parameter inputs to float. In the $E_8$-Persistence theory, every value is derived strictly from the lattice invariants. The Degrees of Freedom (DOF) for this fit is exactly equal to the number of observables ($N=17$).

\begin{equation}
    \chi^2_{total} \approx \flavorTotalChiVal \quad \left( \text{Reduced } \chi^2 \approx \flavorReducedChiVal \right)
\end{equation}

A reduced $\chi^2$ of $\approx \flavorReducedChiVal$ signifies an exceptionally strong fit to empirical reality. The PMNS neutrino sector predictions act cooperatively with the CKM predictions, driving the global statistical confidence closer to unity. 

The primary contribution to the residual arises from the $V_{ud}/V_{cs}$ sector. The geometric framework predicts exact unitarity for the full 4×4 extended matrix (including the sterile channel $\nu R$), but the active 3×3 submatrix exhibits the resolution deficit $1/(\nu\Delta)$ = 1/688. This is the structural origin of the Cabibbo Angle Anomaly: experimental measurements that exclude the sterile channel necessarily find an apparent unitarity violation of precisely this magnitude. If future measurements resolve the anomaly in favor of the 1/688 prediction, this would strongly support the discrete substrate hypothesis.